% !LW recipe = latexmk (xelatex)
\documentclass[12pt,a4paper,UTF8]{ctexart}
\usepackage{SJTUReport}
\usepackage{booktabs} % 三线表
\usepackage{amsmath}
\usepackage{amssymb}

%%-------------------------------正文开始---------------------------%%
\begin{document}

% %-----------------------封面--------------------%%
\cover

%------------------摘要-------------%%
\newpage
\begin{abstract}

面向高速公路长途出行场景，本文研究沿线换电站网络中预约与随机两类异质用户的实时运营优化问题。长途电动汽车在一次行程中可能需要多次换电，每次换电在取走满电电池的同时退回一块低电量电池，使用户的换电站序列选择与沿线各站的电池库存形成时空耦合。为此，本文构建模型预测控制（MPC）与强化学习（RL）相结合的分层协同框架：MPC 在滚动预测域内联合优化预约用户剩余行程的换电站路径与站内连续服务事件安排，RL actor 为各站每个槽位电池生成请求充电功率，critic 则通过分段线性状态价值函数，统一评估期末库存与未完成换电请求对后续运营净收益的影响。模型以连续时间刻画到站、等待、超时、充电完成与瞬时换电事件，允许仍在等待的请求带入下一轮，并按真实发生时刻对服务、违约与路径发布结算一次。

\end{abstract}
\newpage

\thispagestyle{empty} % 首页不显示页码

%%--------------------------目录页------------------------%%
\newpage
\tableofcontents

%%------------------------正文页从这里开始-------------------%
\newpage

%可选择这里也放一个标题
\begin{center}
   \title{ \Huge \textbf{{考虑异质用户的高速公路换电站运营 MPC-RL 分层协同优化}}}
\end{center}


\section{Introduction}
随着电动汽车在高速公路中长距离出行场景中的普及，沿线补能设施的运营能力逐渐成为影响用户出行可靠性与系统运行效率的关键因素。相比传统充电方式，换电模式可以缩短车辆补能时间。在高速公路场景下，长途出行的电动汽车往往需要沿行程多次换电：车辆沿道路单向行驶且不能折返，前一次换电的站点选择直接改变后续可达站点集合与剩余行程的换电站路径，各次换电决策因此相互关联。

与此同时，不同车辆到达时的电池荷电状态（SOC）各不相同：SOC 通过影响剩余续航里程，决定车辆首个可达换电站以及整条换电站序列；而每次换电在换出满电电池的同时退回一块低 SOC 电池，退回电池须经重新充电后才能再次服务。因此，一位用户的站点选择不仅影响自身后续行程，还会改变相关站点未来的电池库存状态，使用户路径与多站点电池库存产生时空耦合。高速公路换电站运营由此成为一个具有用户车辆电池状态和多站点电池状态时空耦合、动态演化特征的序贯决策问题。

实际运营中存在两类异质用户：预约用户和随机到达用户。预约用户在运营日前提交计划到达高速入口的时间、行程起点和终点以及到达 SOC，运营者据此完成接纳决策并发布日前换电路径；随机用户直接到站提出请求，其到达时刻与 SOC 只能预测。预约用户的实际到达信息可能偏离预约内容；车辆进入高速公路后，其未执行的剩余行程的换电站路径需要随实时位置、实时 SOC 和预计到达时间不断更新；到站时若暂无满电电池可用，用户在站内等待，超过其等待耐心后流失。考虑到分时电价使站内充电成本随时间变化，如何协调预约用户履约、随机用户接纳以及站内电池充电调度，是高速公路换电网络运营中的核心问题。

针对上述问题，本文构建 MPC--RL 分层协同优化框架。换电站的实时运营中存在两类性质不同但相互耦合的决策：一类是预约用户剩余换电路径的动态调整决策，即在有限预测域内根据未到用户的预约信息和在途用户的实时状态调整其未执行的换电站序列；另一类是各站每个槽位电池的充电功率决策。剩余路径调整具有离散组合结构，并同时受到车辆续航、路径连通性、预约履约和多站点库存约束；随机需求预测与实际到达又随时间不断更新。模型预测控制（Model Predictive Control, MPC）能够在有限预测域内联合评估这些约束和未来库存影响，并在每次获得新观测后重新优化而只执行当前外层区间，因而适合处理该类动态路径调整问题。充电功率的影响则具有明显的跨时段累积性：当前功率不仅决定当期充电成本，还会改变未来各站每个槽位电池的 SOC、满电电池供给能力和后续服务水平。本文采用强化学习（Reinforcement Learning, RL）学习每个槽位的充电功率控制策略，使其能够在长期交互过程中根据需求、电价和电池状态自适应调整充电行为。如\autoref{fig:MPC-RL}所示，RL actor 输出预测域内的每个槽位的请求充电功率并作为 MPC 的已知参数，MPC 在给定功率下联合优化剩余路径与站内连续服务事件安排，critic 向 MPC 提供预测域之外的长期价值信息。

\begin{figure}[!htbp]
  \centering
  % FIGURE_PROMPT_BEGIN: fig:MPC-RL
  % 绘制一张适用于运筹优化与智能交通学术论文的 MPC–RL 协同框架图。白色背景，扁平化矢量风格，横向 16:9，配色仅使用深蓝、青绿色、橙色和中性灰，无渐变、阴影、3D 或装饰图标。
  % 图中从左到右展示：当前观测状态（每个槽位电池 SOC、预约与随机等待队列及剩余耐心、未到和在途预约信息、实时位置与 ETA、电价和随机需求预测）；基于已知信息预先确定的计划换电路径与计划执行轨迹；RL actor 输出预测域内每个槽位的请求功率，critic 提供固定参数的分段线性价值模型，在 MPC 内评价随决策变化的终端状态；MPC 在给定 RL 信息后优化预约用户剩余路径和站内连续事件安排；确定性执行器在当前外层区间内按真实到站、充电完成、换电和超时事件连续执行；输出实际 reward、更新后的电池状态、等待队列和已发布剩余路径，并反馈至下一轮。
  % 明确区分外层固定间隔决策与区间内部连续事件，强调计划换电路径先确定、随后 actor 滚动推演、最后 MPC 求解，不出现 RL–MPC 内循环。所有文字使用简洁中文，字体统一，箭头方向明确，模块不超过六个，适合缩放后在 A4 论文中阅读。输出优先为可编辑 SVG/PDF 或高分辨率透明背景图片。
  % FIGURE_PROMPT_END: fig:MPC-RL
  \fbox{\parbox[c][0.52\textwidth][c]{0.92\textwidth}{%
    \centering\small Figure placeholder}}
  \caption{MPC--RL 协同决策、连续事件执行与状态反馈框架}
  \label{fig:MPC-RL}
\end{figure}

进一步地，本文利用 RL critic 提供预测时域之外的长期价值信息，以缓解有限预测时域导致的末端效应。传统 MPC 通常通过设置固定终端库存下界避免时域末端过度消耗满电电池，但该方法难以同时反映库存 SOC 分布和未来换电请求的时间差异。本文以预测期末之后的运营净回报训练状态价值函数，统一评估期末库存与未完成请求的长期影响；未来请求保留其到达时间与等待期限，在后续实际服务时才消耗满电库存。为便于嵌入 MPC，critic 采用分段线性结构，其参数在每轮求解中固定，输入的终端状态则随本轮路径与服务决策变化。

因此，本文提出的 MPC--RL 框架基于决策属性进行分层：MPC 负责处理强约束、组合结构明显且需要可解释性的预约用户剩余路径动态调整问题；RL 负责学习具有长期影响的各槽位充电功率策略，并通过 critic 向 MPC 提供后续运营价值信息。二者通过每个槽位 SOC 状态、每个槽位的请求充电功率、在途预约履约事件和终端价值项相互连接，从而在保证预约用户服务质量要求和运营约束可行性的同时，使有限时域 MPC 能够近似感知预测时域之外的长期影响。

% CONTRIBUTIONS_PLACEHOLDER_BEGIN
% 待作者补充本文贡献
% CONTRIBUTIONS_PLACEHOLDER_END

本文其余部分安排如下：第 2 节给出问题描述；第 3 节建立滚动时域下的 MPC 模型，包括基于 SOC 分档的离线候选网络、日前基准路径生成和在线优化模型；第 4 节建立充电控制的马尔可夫决策过程并介绍 RL 算法；第 5 节给出数值实验；第 6 节总结全文。全文将站点规定的满电 SOC 阈值归一化为 $1$；因此，“满电”表示达到该运营目标值，并不要求等于电芯的电化学上限。






\section{Problem Statement}

考虑由多个 O-D 对及沿线换电站构成的单向高速公路换电网络。对于任一 O-D 对 $p$，车辆从入口 $o^p$ 进入高速公路后沿下游方向顺序行驶且不能折返，最终从出口 $d^p$ 离开；一次长途行程可能需要在沿线换电站集合 $\mathcal S^p$ 中的多个站点换电。站点 $i$ 配备槽位集合 $\mathcal B_i$，槽位中的电池以归一化 SOC 刻画电量状态。每次换电从站点取走一块 SOC 为 $1$ 的满电电池，并退回一块带有剩余电量的电池，退回电池入槽充电后可再次服务。因此，某位用户的站点选择不仅决定其后续可达站点，还会改变相关站点未来的电池库存，使用户路径与多站点库存产生时空耦合。

系统服务预约用户和随机到达用户。预约用户在运营日前提交预约信息，包括 O-D 对、预计进入高速公路的时刻及进入时的 SOC，运营者据此决定是否接纳预约，并为已接纳用户制定和发布日前换电路径。随机用户在到达换电站时提出请求，运营者只掌握其预测到达信息。按运行状态，已接纳预约用户分为两类：尚未进入高速公路的未到用户，以及已进入高速公路的在途用户；已经迟于预约入口时刻但尚未实际进入的用户仍属于未到集合，不能仅按原预约时刻将其删除。未到用户的有效入口 SOC 记为 $\widetilde{\mathrm{soc}}_{A,\ell}^{p,k}$，其初值为预约入口 SOC，并在获得新的入口前观测后更新；用户实际进入后，系统改用其实时位置、实时 SOC 和到达下游节点的预计到达时间。到站时若暂无满电电池可用，用户在站内等待：任一请求 $r$ 自其到站时刻 $t_r^{\mathrm{arr}}$ 起最长等待 $\Delta t$，开始服务时刻 $t_r^{\mathrm{sw}}$ 须满足 $t_r^{\mathrm{arr}}\le t_r^{\mathrm{sw}}\le t_r^{\mathrm{ddl}}$，其中 $t_r^{\mathrm{ddl}}=t_r^{\mathrm{arr}}+\Delta t$ 为等待截止时刻；截止时刻仍未获得服务的预约记为一次预约失败，随机请求则流失。换电视为瞬时事件。

系统以固定长度 $\sigma$ 的外层区间组织运营。以运营日统一时间原点，记
\[
t_q=q\sigma,
\qquad
\mathcal I_q=[t_q,t_{q+1}),
\qquad
\mathcal T_\ell=\{\ell,\ldots,\ell+H-1\}.
\]
在每个区间起点 $t_\ell$，系统以未来 $H$ 个区间为预测域，即 $\mathcal T_\ell$ 对应连续时间范围 $[t_\ell,t_{\ell+H})$。预约与随机请求的到站时刻均为外生连续值，区间索引只用于请求归组与电价等分段常数数据。服务顺序遵循两条规则：预约请求优先于随机请求；同类请求内部按原到站时刻升序，同一时刻按固定用户编号升序。系统不为尚未到站的请求扣留当前可用电池；未服务完的等待请求带入下一轮，保留其原到站时刻与原截止时刻，不重新计时。

运营者需要协调预约用户剩余行程的换电站路径与各站每个槽位电池的充电功率，从而在履行已接纳预约的前提下，权衡换电收益、购电成本、路径调整成本和长期库存价值。本文采用分层 MPC--RL 框架：RL 根据需求、电价和库存状态生成预测域内的每个槽位的请求充电功率，并向 MPC 提供分段线性终端价值模型；MPC 在给定请求功率下联合优化剩余路径与站内连续服务事件安排。每轮优化后，系统发布发生变化的在途用户剩余行程的换电站路径，只执行首个外层区间内的充电与服务安排，随后根据真实到站、充电完成、换电和超时事件更新状态并进入下一轮求解。

\autoref{tab:notation}列出正文反复使用的核心符号；仅局部使用的符号在首次出现处就近说明。RL 的状态、动作和算法符号在第 4 节定义。

\begin{table}[!htbp]
  \centering
  \caption{正文主要符号}
  \label{tab:notation}
  \small
  \setlength{\tabcolsep}{3pt}
  \renewcommand{\arraystretch}{0.98}
  \begin{tabular}{@{}lp{11cm}@{}}
    \toprule
    \textbf{符号} & \textbf{含义} \\
    \midrule
    \multicolumn{2}{l}{\textit{集合与索引}} \\
    $\sigma,t_q,\mathcal I_q,\mathcal T_\ell,\Delta t$ & 外层区间长度与起点、半开区间、含 $H$ 个区间的预测域索引集，以及最大等待时间 \\
    $\mathcal P,p;\ \mathcal S^p,\mathcal S;\ \mathcal B_i,b$ & O-D 对、沿线与全系统站点、库存槽位的集合和索引；槽位不是永久电池编号 \\
    $\mathcal K_\ell^{\mathrm{fut},p},\mathcal K_\ell^{\mathrm{enr},p},\mathcal K_\ell^p$ & 未到、在途预约用户集合及其并集；迟到但尚未实际进入者仍属于未到集合 \\
    $o^p,d^p,o_\ell^{\mathrm{cur},p,k},\widetilde{\mathrm{soc}}_{A,\ell}^{p,k},\mathrm{soc}_{A,\ell}^{p,k}$ & 入口、出口、实时虚拟起点、未到用户有效入口 SOC 与在途用户实时 SOC \\
    $\mathcal A^{p,k},t_{A,i}^{p,k},t_r^{\mathrm{arr}},t_r^{\mathrm{ddl}},t_r^{\mathrm{sw}}$ & 个体剩余候选弧集、连续 ETA，以及请求的到站、截止和开始服务时刻 \\
    \midrule
    \multicolumn{2}{l}{\textit{参数与状态}} \\
    $D_{i,j}^{p},v_{i,j}^{p},\underline D,\underline s_d^p$ & 节点间距离与 SOC 消耗、最小换电间距和出口最低 SOC \\
    $e_{i,q},\pi_{i,q}^{\mathrm{sw}},\overline p_i,\overline P_i$ & 购电价、服务价格、单槽功率上限和站级总功率上限 \\
    $E_B,\eta_i,\widehat P_{i,b,q},T_{i,b,q}^{\mathrm{ch}},S_{i,b}(t)$ & 电池额定能量、充电效率、RL 请求功率、槽位在时段内的实际充电时长和槽位 SOC \\
    $\bar{\mathbf y}_{A}^{p,k},\mathbf y_{A}^{\mathrm{pub},p,k}$ & 日前初始发布路径与最近一次实际发布的剩余路径 \\
    $\kappa,c^{\mathrm{fail}},\beta$ & 单次路径调整成本、单次预约失败成本与 RL 终端价值权重 \\
    \midrule
    \multicolumn{2}{l}{\textit{决策变量与主要辅助量}} \\
    $y_{i,j}^{p,k},x_{A,i}^{p,k},d_{p,k}$ & 剩余路径、站点访问量与路径差异指示量 \\
    $s_r,f_r,\omega_r,a_r$ & 请求的服务、失败或流失、带入下一轮三种结局，以及有效服务量 \\
    $s_{\ell+H}^{\mathrm{MPC}},\Phi^{\mathrm{RL}},w_h^{\mathrm{pend}}$ & 本轮决策对应的终端状态、后续运营价值及期末未完成换电指示量 \\
    \bottomrule
  \end{tabular}
\end{table}


\section{MPC Formulation}

设当前决策区间为 $\mathcal I_\ell$。在时刻 $t_\ell$，系统读取：各站每个槽位电池的观测 SOC；自上一轮遗留、仍未超时的站内等待请求；在途预约用户的实时位置、实时 SOC 与到达下游站点的预计到达时间；已发布且尚未执行的预约用户剩余行程的换电路径；以及未来 $H$ 个区间的随机需求预测与电价。

随后，系统按固定顺序构造本轮求解输入。首先确定计划换电路径：在途用户取最近一次实际发布路径；未到用户优先沿用上一轮保留的计划换电路径。其次，以RL actor 的策略，在计划换电路径下滚动生成预测域内的每个槽位的充电功率 $\widehat P_{i,b,q}$。

在此基础上，系统为每位预约用户构造剩余行程的候选路径网络，并由 MPC 更新剩余行程的换电路径安排。

求解后，系统立即发布发生变化的在途用户剩余行程的换电站路径并结算一次路径调整成本，新路径作为下一轮的调整基准；未到用户的优化路径暂不发布。系统只执行首个外层区间 $\mathcal I_\ell$：区间内按真实到站、充电完成、换电和超时事件连续推进；区间结束时，每个槽位 SOC、仍未超时的等待队列与已发布剩余路径进入下一轮。上述过程如\autoref{fig:rolling_horizon}所示。

\begin{figure}[!htbp]
  \centering
  % FIGURE_PROMPT_BEGIN: fig:rolling_horizon
  % 绘制一张适用于学术论文的横向滚动时域示意图。白色背景、简洁矢量风格，宽高比约 2:1，使用深蓝表示预测域、橙色表示当前执行区间、绿色表示带入下一轮的等待请求、灰色表示已执行历史。
  % 顶部为时间轴，清晰标出 $t_{\ell-1}$、$t_\ell$、$t_{\ell+1}$ 和 $t_{\ell+H}$，预测域采用半开区间 $[t_\ell,t_{\ell+H})$，突出只执行首个外层区间。下方设置三条泳道：第一条为未到预约用户，显示其内部路径可滚动优化但仍相对初始发布路径评价；第二条为在途预约用户，显示从实时虚拟起点重优化剩余行程的换电站路径，路径变化后发布并更新比较基准；第三条为站内等待队列，显示原始到站时刻和截止时刻跨滚动边界保持不变。
  % 当前执行区间内部使用少量时间标记展示连续到站、充电完成、立即换电和超时事件；预测终点处显示尚未截止的请求传入下一轮。不得出现“已发布路径固定不再优化”“整时段统一换电”或“随机用户必须当期完成”等旧逻辑。文字精炼、层级清楚、无大段说明，适合 A4 论文单栏宽度阅读。输出优先为可编辑 SVG/PDF 或高分辨率图片。
  % FIGURE_PROMPT_END: fig:rolling_horizon
  \fbox{\parbox[c][0.46\textwidth][c]{0.92\textwidth}{%
    \centering\small Figure placeholder}}
  \caption{滚动时域中的剩余路径优化、连续事件执行与等待请求的跨轮传递}
  \label{fig:rolling_horizon}
\end{figure}


\subsection{Offline Candidate networks by SOC Range}

在每个决策时刻，需要安排预测域内每位预约用户剩余行程的换电路径。如何高效、清晰地表示解的结构并设计决策变量是首要问题。本文借鉴加油站选址模型中 expanded network 的思路，用一条入口--出口路径表示该序列：对于 O-D 对 $p$，若路径为
$o^p\rightarrow i_1\rightarrow\cdots\rightarrow i_r\rightarrow d^p$，则用户依次在换电站 $i_1,\ldots,i_r$ 换电，一条弧跨过的沿线站点均不被访问。为避免在每个滚动时刻重复枚举全部下游组合，本文在每轮优化前先生成候选路径网络。不同的用户出发 SOC 会影响车辆的续航里程，从而影响可达的换电站集合。我们将用户出发 SOC 划分为若干个离线区间，并在每个区间上生成相应的候选网络。

对每个 O-D 对 $p$，节点集合为
$\mathcal V^p=\{o^p\}\cup\mathcal S^p\cup\{d^p\}$，所有节点均按高速公路下游方向排列。参数 $D_{i,j}^p$ 和 $v_{i,j}^p>0$ 分别表示从 $i$ 到下游节点 $j$ 的行驶距离和 SOC 消耗。

考虑到现实场景中为用户安排两个距离过近的连续站点换电并不合理，同时也为了减少网络规模，给定最小换电间距 $\underline D>0$，生成的网络在保证方案可行性的基础上，尽量排除从出发位置到首个换电站以及任意两个连续访问换电站之间距离小于 $\underline D$ 的方案。指向出口的末段行程不受该间距约束，但车辆到达出口时的 SOC 不得低于 $\underline s_d^p$。

对每个 SOC 区间 $h$ 和 O-D 对 $p$，离线网络按以下步骤生成。

\begin{enumerate}
  \item[\textbf{Step 1.}] \textbf{生成原始弧。}若车辆以该 SOC 区间的上限电量能够从入口 $o^p$ 到达下游换电站 $j$，则连接二者；若 $v_{i,j}^p\le1$，则连接两个依次位于下游的换电站 $i$ 和 $j$；若 $v_{i,d^p}^p+\underline s_d^p\le1$，则连接换电站 $i$ 与出口 $d^p$。

  \item[\textbf{Step 2.}] \textbf{生成完整可行路径。}在原始网络中枚举所有从 $o^p$ 到 $d^p$ 的完整路径。

  \item[\textbf{Step 3.}] \textbf{剪枝过近换电弧。}从上述完整路径中找出距离小于 $\underline D$ 的首段弧和站间弧，如果现有的完整路径中仍存在至少一条不包含该弧的方案，则删除该弧，否则保留。

  \item[\textbf{Step 4.}] \textbf{形成候选网络。}将最终保留的完整路径所包含的弧合并为 $\mathcal A^{h,p}$，候选网络由节点集 $\mathcal V^p$ 和弧集 $\mathcal A^{h,p}$ 组成。
\end{enumerate}

离线生成的候选网络刻画了站点之间的下游连接关系，对尚未进入高速的用户，直接按其有效入口 SOC $\widetilde{\mathrm{soc}}_{A,\ell}^{p,k}$ 所在的档位选取对应的离线网络。对已经在途的用户，将当前位置设置为虚拟入口，将当前 SOC 看作有效入口SOC。针对每位预约用户选取的候选弧集记为 $\mathcal A^{p,k}$。


\subsection{Day-Ahead Baseline Path Generation}

用户提交预约信息
$\left(p,\bar t_A^{p,k},\overline{\mathrm{soc}}_A^{p,k}\right)$
后，运营者依据其行程、预计入口时刻和入口 SOC 检查能否提供服务，并为可接纳的请求生成日前基准路径。

预约请求按提交顺序处理。处理当前请求前，算法根据已接纳预约和日前随机需求预测，递推估计各站未来各时段的电池库存。日前充电功率由 RL 在日前预测状态上逐区间生成。电池未满时按当前区间的日前功率充电，达到 SOC $1$ 后立即停止充电，换电退回的低 SOC 电池入槽后立即恢复充电。

对 O-D 对 $p$ 的预约用户 $k$，基准路径按以下规则生成：
\begin{enumerate}
  \item 车辆从入口出发时的 SOC 为 $\overline{\mathrm{soc}}_A^{p,k}$，每次换电后恢复为 $1$。下一节点沿候选弧集选择，因此首段和站间行程均满足续航要求。
  \item 若车辆能够从当前位置直接到达出口，且到达 SOC 不低于 $\underline s_d^p$，则下一站选择出口；否则，仅保留预计到站时刻已经有满电电池可用的候选站。
  \item 候选站按预计满电电池余量从高到低排序，余量相同时优先选择更靠近出口的站点。选定站点后，在相应时段为该用户预留一块满电电池，并更新该站后续时段的预计库存；随后重复上述过程。
\end{enumerate}

若任一步不存在满足库存与下游可达性要求的候选站，即不能生成完整入口--出口路径，则拒绝该请求；只有成功生成完整路径的用户才进入已接纳集合。对已接纳用户，若基准路径经过弧 $(i,j)$，则令 $\bar y_{i,j}^{p,k}=1$，否则为 $0$。该路径随预约结果一并发布，既作为用户的初始发布路径，也作为其未到阶段路径调整的固定比较基准。


\subsection{Optimization Model}


\paragraph{决策变量。}
本轮 MPC 优化的核心决策为预约用户剩余行程的换电路径。决策变量表示为
\[
y_{i,j}^{p,k}\in\{0,1\},
\qquad
(i,j)\in\mathcal A^{p,k},
\quad
k\in\mathcal K_\ell^p.
\]
若 O-D 对 $p$ 的预约用户 $k$ 的剩余行程的换电路径经过弧 $(i,j)$，则 $y_{i,j}^{p,k}=1$；一组满足流平衡的 $y$ 变量唯一确定该用户依次访问的换电站序列。

服务请求用索引 $r$ 表示，包括预测域内预计到达的预约与随机请求，以及上一轮遗留、仍未超时的等待请求。每个请求分别用以下二元变量表示：$s_r$ 表示在预测域内获得服务，$f_r$ 表示超过等待耐心而失败（预约请求）或流失（随机请求），$\omega_r$ 表示到预测期末仍在等待、带入下一轮；请求的开始服务时刻记为 $t_r^{\mathrm{sw}}$。此外，$d_{p,k}$ 标记用户 $k$ 剩余行程的换电站路径相对其比较基准是否发生变化。各槽位电池的电量记为 $S_{i,b}(t)$。

每个槽位的请求充电功率 $\widehat P_{i,b,q}$ 与分段线性 critic 的参数由 RL 给出，在本轮 MPC 求解中保持固定；终端状态及其价值随决策变化。电池未满时按当前时段的请求功率线性充电，充满即停；换入退回电池后立即恢复充电。

\paragraph{目标函数。}
本轮 MPC 的目标由服务收益、充电成本、路径调整成本、预约超时成本和 RL 终端价值五项组成。

\textbf{换电服务收益。}每完成一次换电，站点按服务时刻所在时段的服务价格 $\pi_{i,q(t_r^{\mathrm{sw}})}$ 对换出电量收费：用户取走满电电池、退回电量为 $\rho_r$ 的电池时，本次服务电量为 $E_B(1-\rho_r)$。预测域内的服务收益为
\begin{equation}
I
=E_B\sum_{i\in\mathcal S}
\sum_{r\in\mathcal E_i}
\pi_{i,q(t_r^{\mathrm{sw}})}^{\mathrm{sw}}(1-\rho_r)\,s_r,
\label{eq:income}
\end{equation}
其中 $\mathcal E_i$ 为站点 $i$ 的候选请求集合，$q(t)$ 表示时刻 $t$ 所属的时段。

\textbf{充电成本。}记 $T_{i,b,q}^{\mathrm{ch}}$ 为槽位 $b$ 在时段 $q$ 内的实际充电时长（扣除已经满电、停止充电的时间），站点 $i$ 在时段 $q$ 的电价为 $e_{i,q}$，则预测域内的购电成本为
\begin{equation}
C_{\mathrm{ch}}
=
\sum_{i\in\mathcal S}
\sum_{q\in\mathcal T_\ell}
e_{i,q}
\sum_{b\in\mathcal B_i}
\widehat P_{i,b,q}\,T_{i,b,q}^{\mathrm{ch}}.
\label{eq:chargingcost}
\end{equation}


\textbf{路径调整成本。}若预约用户剩余行程的换电路径不同于上一次运营商为其发布的行程，则产生一次调整成本：
\begin{equation}
C_{\mathrm{adj}}
=
\kappa
\sum_{p\in\mathcal P}
\sum_{k\in\mathcal K_\ell^p}
d_{p,k}.
\label{eq:adjustment_cost}
\end{equation}
未到用户的比较基准是日前发布路径，在途用户的比较基准是最近一次实际发布路径。

\textbf{预约超时成本。}当预约请求在其等待截止时刻仍未获得服务时，记一次预约失败：
\begin{equation}
C_{\mathrm{fail}}
=
c^{\mathrm{fail}}
\sum_{i\in\mathcal S}
\sum_{r\in\mathcal E_{i}^{A}}
f_r,
\label{eq:reservation_failure_cost}
\end{equation}
其中 $\mathcal E_i^A$ 为站点 $i$ 的预约请求子集。

\textbf{RL 终端价值。}有限预测域截断了两类长期影响：期末剩余库存对后续运营的影响，以及期末未完成换电请求在未来服务时对库存和收益的影响。本文采用第 4 节训练的状态价值函数 $V_\theta$，通过期末之后的运营净回报统一衡量这两类影响。

记本轮路径与服务决策对应的预测终端状态为 $s_{\ell+H}^{\mathrm{MPC}}$，其包含期末电池 SOC、未完成请求及其到达时间、截止时间和退回 SOC，以及期末时间、电价和后续需求预测。该状态沿用第 4 节的状态信息与编码口径，但其中受本轮决策影响的库存、请求状态和剩余路径均由 MPC 变量确定，不固定为计划路径下的参考推演结果。未完成请求在期末仍作为待服务请求保留，不提前扣除满电库存，也不提前写入退回电池；其后续服务、等待与超时的影响由价值函数统一评价。

进入目标的终端价值为
\begin{equation}
\Phi^{\mathrm{RL}}
=V_\theta\!\left(s_{\ell+H}^{\mathrm{MPC}}\right).
\label{eq:rl_terminal_value}
\end{equation}
$V_\theta$ 采用式\eqref{eq:pwl_critic}的分段线性结构，并通过式\eqref{eq:pwl_critic_embedding}嵌入求解。期末之后的服务收入、充电与路径调整费用及预约超时损失由该价值估计，域内已结算项目不重复计入。

为确定终端状态中的未完成预约请求，令 $\overline{\mathcal D}_{\ell}^{A}$ 为由本轮候选弧和遗留预约请求预先枚举的候选换电集合，$w_h^{\mathrm{pend}}\in\{0,1\}$ 表示其中第 $h$ 项到预测期末仍有效且未完成服务。严格地，令 $A_{p,k}^{\mathrm{bd}}\in\{0,1\}$ 表示用户 $(p,k)$ 在预测域内未发生超时：
\begin{subequations}
\label{eq:boundary_alive}
\begin{align}
A_{p,k}^{\mathrm{bd}}
&\ge 1-\sum_{\varepsilon\in\mathcal E_{p,k}^{A,\mathrm{in}}}
f_\varepsilon,\\
A_{p,k}^{\mathrm{bd}}
&\le 1-f_\varepsilon,
\qquad \varepsilon\in\mathcal E_{p,k}^{A,\mathrm{in}},
\end{align}
\end{subequations}
其中 $\mathcal E_{p,k}^{A,\mathrm{in}}$ 为该用户在预测期末前可能到达的候选事件集合，空集时固定为 $1$。对 $h\in\overline{\mathcal D}_{\ell}^{A}$，令 $\delta_h$ 表示该次换电是否被路径选中（普通候选事件取入弧变量 $y$，遗留请求取 $1$），$s_h^{\mathrm{in}}$ 为其域内服务量（期末前不可能服务的固定为 $0$），则
\begin{subequations}
\label{eq:pending_delivery_alive}
\begin{align}
w_h^{\mathrm{pend}}&\le\delta_h,&
w_h^{\mathrm{pend}}&\le A_{p(h),k(h)}^{\mathrm{bd}},\\
w_h^{\mathrm{pend}}&\le1-s_h^{\mathrm{in}},&
w_h^{\mathrm{pend}}
&\ge\delta_h+A_{p(h),k(h)}^{\mathrm{bd}}
-s_h^{\mathrm{in}}-1.
\end{align}
\end{subequations}
即期末未完成换电是“路径选中、用户未失效、域内尚未服务”三者的交集。预约请求按 $w_h^{\mathrm{pend}}$ 纳入终端状态，随机等待请求按 $\omega_r$ 纳入；已到站的预约等待请求只计一次，同一用户尚未完成的下游换电保留其顺序与上游履约依赖。$\beta\ge0$ 为终端价值权重；域内金额保持未折扣口径，因此本目标是加权近似运营目标，而非统一折扣回报的严格时域截断。

综合上述各项，本轮 MPC 的目标为最大化综合运营价值：
\begin{equation}
\max\quad
I-C_{\mathrm{ch}}-C_{\mathrm{adj}}-C_{\mathrm{fail}}
+\beta\Phi^{\mathrm{RL}}.
\label{eq:objective}
\end{equation}


站内服务遵循固定规则：预约请求优先于随机请求；同类内部按原到站时刻升序，同时刻按固定编号升序。每次服务使用编号最小的满电槽位。只要还有已到站、未超时的等待请求且有空闲的满电槽位，就必须立即分配服务。换电为瞬时事件，服务时刻 $t_r^{\mathrm{sw}}$ 由系统根据当前状态和规则确定。

\paragraph{约束条件。}

 
（1）\textbf{网路流平衡约束。} 每位预约用户在其候选网络上满足流平衡约束：
\begin{equation}
\sum_{\{j\mid(n,j)\in\mathcal A^{p,k}\}}
y_{n,j}^{p,k}
-
\sum_{\{j\mid(j,n)\in\mathcal A^{p,k}\}}
y_{j,n}^{p,k}
=
\begin{cases}
1,&n=o_\ell^{p,k},\\
-1,&n=d^p,\\
0,&\text{其余节点},
\end{cases}
\quad
p\in\mathcal P,\ k\in\mathcal K_\ell^p,\ n\in\mathcal V^{p,k},
\label{eq:flow}
\end{equation}
其中本轮起点 $o_\ell^{p,k}$ 对未到用户为入口 $o^p$，对在途用户为实时虚拟起点 $o_\ell^{\mathrm{cur},p,k}$；$\mathcal V^{p,k}$ 为该用户剩余候选网络的节点集。退回电池的电量按“出发电量减去该段耗电”计算：用户 $k$ 沿入弧 $(j,i)$ 到达站点 $i$ 时，
\begin{equation}
\rho_{A,j,i}^{p,k}
=
\begin{cases}
\widetilde{\mathrm{soc}}_{A,\ell}^{p,k}-v_{o^p,i}^{p},
&j=o^p,\ k\in\mathcal K_\ell^{\mathrm{fut},p},\\[2pt]
\mathrm{soc}_{A,\ell}^{p,k}
-v_{o_\ell^{\mathrm{cur},p,k},i}^{p,k},
&j=o_\ell^{\mathrm{cur},p,k},\ k\in\mathcal K_\ell^{\mathrm{enr},p},\\[2pt]
1-v_{j,i}^{p},
&j\in\mathcal S^p,\ k\in\mathcal K_\ell^{p},\\[2pt]
\end{cases}
\quad (j,i)\in\mathcal A^{p,k}.
\label{eq:return_soc}
\end{equation}
即首段用用户当前电量，之后每段从上一站换得的满电出发。

（2）\textbf{路径调整约束。}定义当前路径的站点访问量
\begin{equation}
x_{A,i}^{p,k}
=
\sum_{\{j\mid(j,i)\in\mathcal A^{p,k}\}}y_{j,i}^{p,k},
\qquad i\in\mathcal S^p,
\label{eq:station_visit_indicator}
\end{equation}
以及基准访问量
\begin{equation}
x_{A,i\mid\ell}^{\mathrm{ref},p,k}
=
\begin{cases}
\bar x_{A,i}^{p,k},
&k\in\mathcal K_\ell^{\mathrm{fut},p},\\
x_{A,i\mid\ell}^{\mathrm{pub},p,k},
&k\in\mathcal K_\ell^{\mathrm{enr},p},
\end{cases}
\label{eq:reference_visit_indicator}
\end{equation}
其中 $\bar x_{A,i}^{p,k}$ 由日前发布路径汇总，$x_{A,i\mid\ell}^{\mathrm{pub},p,k}$ 由最近一次实际发布路径的未执行站序汇总。路径差异标记满足
\begin{subequations}
\label{eq:path_adjustment_indicator}
\begin{align}
d_{p,k}
&\ge x_{A,i}^{p,k}-x_{A,i\mid\ell}^{\mathrm{ref},p,k},
&&p\in\mathcal P,\ k\in\mathcal K_\ell^p,\ i\in\mathcal S^p,\\
d_{p,k}
&\ge x_{A,i\mid\ell}^{\mathrm{ref},p,k}-x_{A,i}^{p,k},
&&p\in\mathcal P,\ k\in\mathcal K_\ell^p,\ i\in\mathcal S^p.
\end{align}
\end{subequations}
只要剩余行程的换电站路径相对基准增加或删除任意换电站，就有 $d_{p,k}=1$；与基准一致时，在目标函数中的正调整成本作用下有 $d_{p,k}=0$。

（3）\textbf{服务窗口约束。}任一请求自到站时刻起最长等待 $\Delta t$，截止时刻为
\begin{equation}
t_r^{\mathrm{ddl}}
=
t_r^{\mathrm{arr}}+\Delta t,
\label{eq:waiting_deadline}
\end{equation}
获得服务的请求满足 $t_r^{\mathrm{arr}}\le t_r^{\mathrm{sw}}\le t_r^{\mathrm{ddl}}$。三种结局互斥且穷尽：
\begin{equation}
s_r+f_r+\omega_r=a_r,
\label{eq:reservation_conservation_main}
\end{equation}
其中 $a_r$ 为有效服务量，表示该候选请求确实发生。站点 $i$ 的候选请求集合 $\mathcal E_i$ 在求解前由已知信息枚举，包括用户候选网络中指向 $i$ 的弧对应的预约换电事件、预测随机请求和上一轮遗留的等待请求三部分，不以路径变量为成员条件。对随机请求和遗留等待请求，$a_r=1$；对预约事件，首站事件（入弧起点为本轮起点 $o_\ell^{p,k}$）只取决于路径选择，$a_\varepsilon=y_{j,i}^{p,k}$，其余沿弧 $(j,i)$ 从上游站 $j$ 到达的事件满足
\begin{subequations}
\label{eq:reservation_alive_chain}
\begin{align}
a_\varepsilon&\le y_{j,i}^{p,k},\\
a_\varepsilon&\le
\sum_{\varepsilon'\in\mathcal E^A(p,k,j)}s_{\varepsilon'},\\
a_\varepsilon&\ge y_{j,i}^{p,k}
+\sum_{\varepsilon'\in\mathcal E^A(p,k,j)}s_{\varepsilon'}-1,
\end{align}
\end{subequations}
其中 $\mathcal E^A(p,k,j)$ 为该用户在上游站 $j$ 的候选事件集合。这组约束的含义是：用户只有在上游站成功换电，下游事件才有效；一旦上游超时，下游事件全部失效，失败只结算一次；无效事件满足 $s_\varepsilon=f_\varepsilon=\omega_\varepsilon=0$。等待窗口跨预测期末的请求，本轮不能被判定失败或流失，只能在域内获得服务或以 $\omega_r=1$ 带入下一轮。

（4）\textbf{充电与换电约束。}在相邻事件时刻 $t_1<t_2$ 之间（同属于某个时段 $q$），槽位电量线性上升、充满即止：
\begin{equation}
S_{i,b}(t_2)
=
\min\left\{
1,\;
S_{i,b}(t_1)
+
\frac{\eta_i\widehat P_{i,b,q}}{E_B}(t_2-t_1)
\right\},
\label{eq:event_charging_transition}
\end{equation}
即电池未满时按当前时段的功率充电，充满即停。功率的可行范围为
\begin{equation}
0\le \widehat P_{i,b,q}\le\overline p_i,
\qquad
\sum_{b\in\mathcal B_i}\widehat P_{i,b,q}
\le\overline P_i,
\qquad i\in\mathcal S,\ q\in\mathcal T_\ell.
\label{eq:power_projection}
\end{equation}
只有电量恰好为 $1$ 的电池可以交付。换电为瞬时事件：服务完成后槽位电量立即变为退回值，$S_{i,b}\big((t_r^{\mathrm{sw}})^+\big)=\rho_r$，并按当前时段的请求功率恢复充电。槽位是存放电池的位置而非固定的电池，同一槽位在一个决策区间内可多次服务，但两次服务之间必须重新充满。

本轮各槽位的初始电量、正在充电的状态和遗留等待队列均取真实观测。终端状态 $s_{\ell+H}^{\mathrm{MPC}}$ 取预测期末事件发生前的槽位电量 $S_{i,b}(t_{\ell+H}^{-})$、仍在等待及未来尚未到站的有效请求，并保留相应的时间和路径信息；期末未完成预约换电的界定见式\eqref{eq:pending_delivery_alive}。



\section{RL Formulation}

RL 层用于学习每个槽位的充电功率控制策略，并通过分段线性 critic 为 MPC 提供期末之后的运营价值。本节将充电功率控制问题建模为马尔可夫决策过程，并给出本文采用的 RL 算法。

\subsection{Markov decision process}

\paragraph{State.}
在滚动决策时刻 $t_\ell$，RL agent 观测系统状态 $s_\ell$。除各站每个槽位电池的 SOC、电价和时间特征外，状态包含：未到与在途预约用户信息（在途用户含实时位置、实时 SOC 和最近发布的剩余行程的换电站路径）、站内预约与随机等待队列（含原到站时刻与剩余耐心）、最近一个外层区间的真实随机到达记录，以及未来逐请求随机需求预测。令
$\mathcal K_\ell^{\mathrm{rem},p}$ 表示时刻 $t_\ell$ 尚未进入高速公路的已接纳预约用户集合（覆盖范围超出当前预测域），用 $\bar{\mathbf y}_A^{p,k}$ 和 $\mathbf y_A^{\mathrm{pub},p,k}$ 分别汇集日前初始发布路径与最近实际发布路径的弧指示量，则预约信息集合为
\begin{equation}
\begin{aligned}
\mathcal Z_{A}
={}&
\left\{
\left(p,k,\bar t_A^{p,k},
\overline{\mathrm{soc}}_A^{p,k},
\bar{\mathbf y}_A^{p,k}\right):
k\in\mathcal K_\ell^{\mathrm{rem},p}
\right\}\\
&\cup
\left\{
\left(p,k,o_\ell^{\mathrm{cur},p,k},\mathrm{soc}_{A,\ell}^{p,k},
\mathbf y_A^{\mathrm{pub},p,k}\right):
k\in\mathcal K_\ell^{\mathrm{enr},p}
\right\}.
\end{aligned}
\label{eq:rl_reservation_information}
\end{equation}
站内等待队列记为
$\mathcal W_\ell=\{(r,i,t_r^{\mathrm{arr}},t_r^{\mathrm{ddl}})\}$，
其中元素为仍未服务且未超时的预约与随机请求。由于 $\mathcal Z_{A}$ 与 $\mathcal W_\ell$ 的元素数量随滚动时刻变化，分别使用 $f_A$、$f_W$ 将其编码为固定维度特征。再令 $f_B$、$f_R^{\mathrm{act}}$ 和 $f_R^{\mathrm{pred}}$ 分别编码每个槽位 SOC、真实随机到达记录和预测随机请求，则状态表示为
\begin{equation}
\begin{aligned}
s_\ell=\Big(&
f_B\!\left(\{S_{i,b}^{\mathrm{obs}}(t_\ell)\}_{i\in\mathcal S,b\in\mathcal B_i}\right),
f_A(\mathcal Z_{A}),
f_W(\mathcal W_\ell),\\
&
f_R^{\mathrm{act}}\!\left(
\{(i,r,t_R^{i,r},\mathrm{soc}_R^{i,r},z_{R,r,i}^{\mathrm{act}}):
r\in\mathcal R_{i,\ell-1}^{\mathrm{act}}\}_{i\in\mathcal S}
\right),\\
&
f_R^{\mathrm{pred}}\!\left(
\{(i,\nu,t_{R}^{i,\nu},\mathrm{soc}_{R}^{i,\nu}):
\nu\in\mathcal R_{i,q}\}_{i\in\mathcal S,q\in\mathcal T_\ell}
\right),
\{e_{i,q}\}_{i\in\mathcal S,q\in\mathcal T_\ell},
\xi_\ell
\Big).
\end{aligned}
\label{eq:rl_state}
\end{equation}
其中 $\xi_\ell$ 表示时段、节假日或交通高峰等外生特征；$\mathcal R_{i,\ell-1}^{\mathrm{act}}$ 为上一区间真实到站的随机请求记录，$z_{R,r,i}^{\mathrm{act}}$ 为其实际服务结果。$f_R^{\mathrm{act}}$ 使用真实的到达时刻、SOC 和服务结果，$f_R^{\mathrm{pred}}$ 仅使用时刻 $t_\ell$ 可得的预测量，两者不共享数据字段。由于同一站点的槽位标签没有物理意义，$f_B$ 对各槽位 SOC 使用共享的分段线性基函数并在站内求和，以刻画 SOC 分布；actor 的每个槽位的输出仍采用相应的置换等变解码，滚动推演按预先固定的对称性消除规则将事件分配给槽位。各状态编码采用仿射变换、ReLU 与求和构成，使 critic 从原始状态到价值输出保持分段线性。请求编码保留相对于当前状态时刻的到达时间与截止时间、退回 SOC、用户类别及剩余路径关系；同一次换电请求在预约信息与等待队列中只计一次。在 MPC 中，候选请求的固定属性可预先编码，再由 $w_h^{\mathrm{pend}}$、$\omega_r$ 及路径变量控制其是否计入；随决策变化的库存和名义服务记录仍需同步编码。训练与终端评价采用相同的特征定义和时间基准，后续需求预测覆盖期末之后，不能把所有未完成请求合并为无时间信息的总量。

\paragraph{Action.}
RL 的动作是各站点、各槽位在当前外层区间的连续请求充电功率：
\begin{equation}
a_\ell=
\{\widehat P_{i,b,\ell}\}_{i\in\mathcal S,b\in\mathcal B_i}.
\label{eq:rl_action}
\end{equation}
actor 先生成每个槽位的无约束请求，再将其欧氏投影到式\eqref{eq:power_projection}所定义的闭凸集上，并以该唯一投影作为动作 $\widehat P$；该确定映射同时满足单槽功率上限和站级总功率上限。策略概率密度及 PPO 比率均针对这一固定投影参数化所诱导的请求动作分布计算。

\paragraph{Rollout.}
为向 MPC 提供未来 $H$ 个区间的请求功率，系统从
$s_{\ell}^{\mathrm{ref}}=s_\ell$ 出发，利用同一 actor 的均值动作和推演状态递推进行滚动推演：
\begin{equation}
\{\widehat P_{i,b,q}\}_{i,b}
=
\mu_\phi(s_{q}^{\mathrm{ref}}),
\quad q\in\mathcal T_\ell,
\label{eq:rl_power_rollout}
\end{equation}
其中 $\mu_\phi$ 表示策略分布的均值。推演遵循固定的调用顺序：先按第 3 节规则确定计划换电路径，再进行 actor 滚动推演，最后将固定功率序列与固定参数的 critic 一并交给 MPC。计划换电路径的确定只使用求解前已知的位置、SOC、预计到达时间和候选网络；推演采用与实际执行相同的预约优先、等待、超时、充电与电池更换规则，并允许预约软失败。因此，推演状态 $s_q^{\mathrm{ref}}$ 与本轮待求的路径和服务决策无关，请求序列生成与本轮路径优化之间不存在循环依赖。构造后续推演状态时，区间 $q$ 的预测请求及其名义服务结果进入独立的名义历史编码器 $f_R^{\mathrm{nom}}$；该编码器与 $f_R^{\mathrm{act}}$ 输出维度相同，但不共享输入字段或语义，且只用于推演状态和预测终端状态，不写回真实状态。$f_R^{\mathrm{pred}}$ 所需的后续预测窗口由同一需求预测器的条件均值补齐，超出当前 $H$ 步输入的部分采用固定终端预测先验。整条请求功率序列作为本轮 MPC 参数；只有首项请求 $\widehat P_{i,b,\ell}$ 进入真实执行，下一滚动时刻重新生成序列。

\paragraph{Transition.}
给定 $s_\ell$ 和请求功率序列后，系统求解 MPC，得到逐用户最优路径
$\{y_{i,j}^{p,k*}\}$。在途用户中剩余行程的换电站路径发生变化者在本轮边界立即发布并结算一次路径调整成本，新路径作为下一轮推演与调整的基准；未到用户的路径只用于当前预测，不对用户发布。为避免内部预测路径进入真实执行，定义边界后的执行路径：在途用户使用本轮已经发布的最优剩余行程的换电站路径，未到用户仍使用日前发布路径 $\bar{\mathbf y}_{A}^{p,k}$；若后者在 $\mathcal I_\ell$ 内实际进入高速公路，仍先按该已发布路径执行，下一边界再作为在途用户更新。记两类用户在 $\mathcal I_\ell$ 内的真实到站记录分别为 $\mathcal H_{A,\ell}^{\mathrm{act}}$ 和 $\mathcal H_{R,\ell}^{\mathrm{act}}$，每条记录携带请求标识、站点、真实到站时刻和真实退回 SOC。

环境随后只在首个外层区间 $\mathcal I_\ell$ 内执行。确定性连续事件算子 $\operatorname{Exec}_\ell$ 的输入仅为真实到站记录、遗留等待队列、边界后已发布的执行路径、当前请求功率和观测 SOC；其内部按充电完成、登记到站、队列分配、超时判定的同时刻顺序连续推进：
\begin{equation}
\begin{aligned}
(\mathbf S^{\mathrm{obs}}(t_{\ell+1}),\mathcal W_{\ell+1},
\mathbf s_{A,\ell}^{\mathrm{act}},\mathbf z_{R,\ell}^{\mathrm{act}},
\mathbf f_{A,\ell}^{\mathrm{act}})
=\operatorname{Exec}_\ell\!\Big(&
\mathbf S^{\mathrm{obs}}(t_\ell),\mathcal W_\ell,
\{\widehat P_{i,b,\ell}\}_{i,b},\\
&\{\mathbf y_{A,\ell}^{\mathrm{exec},p,k}\}_{p,k},
\mathcal H_{A,\ell}^{\mathrm{act}},\mathcal H_{R,\ell}^{\mathrm{act}}
\Big).
\end{aligned}
\label{eq:actual_execution_operator}
\end{equation}
其中 $\mathbf y_{A,\ell}^{\mathrm{exec},p,k}=\mathbf y_A^{p,k*}$ 对在途用户成立，而对未到用户固定为 $\bar{\mathbf y}_{A}^{p,k}$。实际执行事件只来自两类真实到站记录和遗留等待队列；预测集合不作为真实到达输入执行算子，预测偏差不触发停止执行，也不要求预测与实际终态强制相等。状态转移可表示为
\begin{equation}
s_{\ell+1}
\sim
\mathbb P\!\left(
\cdot\mid
s_\ell,a_\ell,
\mathrm{MPC}(s_\ell,\widehat{\mathbf P})
\right),
\label{eq:rl_transition}
\end{equation}
其中 $\mathbb P$ 表示随机到达与下一轮预测更新共同诱导的环境状态转移核，$\widehat{\mathbf P}=\{\widehat P_{i,b,q}:i\in\mathcal S,b\in\mathcal B_i,q\in\mathcal T_\ell\}$。这一定义与 MPC 的“优化完整预测域、只执行当前区间”保持一致。

\paragraph{Reward.}
RL reward 使用当前实际执行区间的已实现运营收益，并计入本轮实际发布路径产生的调整成本：
\begin{equation}
r_\ell
=
I_\ell^{\mathrm{act}}
-C_{\mathrm{ch},\ell}
-C_{\mathrm{adj},\ell}
-C_{\mathrm{fail},\ell}.
\label{eq:rl_reward}
\end{equation}
其中 $I_\ell^{\mathrm{act}}$、$C_{\mathrm{ch},\ell}$、$C_{\mathrm{adj},\ell}$、$C_{\mathrm{fail},\ell}$ 分别为区间 $\mathcal I_\ell$ 内真实发生的服务收益、充电成本、路径调整成本和预约超时成本，各项均按真实发生时刻结算，且跨轮只结算一次：服务收益按实际服务时刻所在时段的价格结算，充电成本按区间内的实际充电量结算，路径调整成本只在真实向在途用户发布新路径时结算，预约失败成本按截止时刻所在的区间结算；上一轮到站、本轮才服务的请求只在服务发生的区间计一次收入，到预测期末仍在等待的请求不作最终统计。MPC 目标中的预测随机收益、未来预约超时风险和未来区间成本不进入真实 reward；未到用户的前瞻调整成本也不在当期 reward 中重复实现。该 reward 使 actor 学习每个槽位的充电功率对当前成本、预约可靠性、路径稳定性和长期电池状态价值的共同影响。


\subsection{RL algorithm}

由于每个槽位的请求功率是连续动作，且需要在不确定需求环境下保持训练稳定性，本文采用 Proximal Policy Optimization (PPO) 作为 RL 层的训练算法。PPO 属于 actor-critic 方法，actor $\pi_{\phi}(a_t\mid s_t)$ 输出连续请求功率动作的随机策略，critic $V_{\theta}(s_t)$ 估计当前状态的长期价值。其 critic 直接学习 state value function，可由式\eqref{eq:rl_terminal_value}统一评价终端库存与未完成请求的后续影响。该价值对应 actor 与 MPC 组成的后续闭环控制下的期望折扣运营净回报，不另设逐槽位价值或未完成请求的独立扣减项。

为直接嵌入 MPC，critic 采用具有 $M$ 个 ReLU 单元的分段线性模型：
\begin{equation}
V_\theta(s)
=\boldsymbol a^\top s+b
+\sum_{m=1}^{M}c_m
\max\{0,\boldsymbol w_m^\top s+d_m\}.
\label{eq:pwl_critic}
\end{equation}
其中 $\theta$ 汇集模型参数，每个单元联合读取库存、请求与时间特征，以表达它们之间的相互作用，不预设价值函数为凸或凹。在一次 MPC 求解中固定全部网络及编码参数，令 $\zeta_m=\boldsymbol w_m^\top s_{\ell+H}^{\mathrm{MPC}}+d_m$，以 $u_m$ 表示其 ReLU 输出。给定有效界 $L_m\le\zeta_m\le U_m$，当 $L_m<0<U_m$ 时，引入 $\delta_m^V\in\{0,1\}$，有
\begin{subequations}
\label{eq:pwl_critic_embedding}
\begin{align}
u_m&\ge0,& u_m&\ge\zeta_m,\\
u_m&\le U_m\delta_m^V,&
u_m&\le\zeta_m-L_m(1-\delta_m^V),
\qquad m=1,\ldots,M.
\end{align}
\end{subequations}
其中 $L_m,U_m$ 由本轮可行状态的有效上下界推得；若 $U_m\le0$，直接取 $u_m=0$，若 $L_m\ge0$，直接取 $u_m=\zeta_m$。将式\eqref{eq:pwl_critic}中的最大值替换为 $u_m$，并对状态编码中的 ReLU 作同样处理，即可用混合整数线性约束表示终端价值模块；其规模随编码及网络中需判定激活状态的单元数增加。

PPO 的 critic 通过最小化价值函数误差进行训练：
\begin{equation}
L_V(\theta)=
\mathbb{E}_t\left[
\left(V_{\theta}(s_t)-\hat{G}_t\right)^2
\right],
\label{eq:value_loss}
\end{equation}
其中 $\hat{G}_t$ 为由采样轨迹中后续实际运营 reward 计算的折扣回报目标，服务收入、充电费用及超时损失均按实际发生时刻计入。采样批次内固定供 MPC 使用的网络参数，批次结束后更新模型；训练样本需覆盖不同库存分布、请求到达时间与等待压力。actor 采用 clipped surrogate objective 更新策略。定义重要性采样比率
\begin{equation}
w_t^{\mathrm{IS}}(\phi)=
\frac{\pi_{\phi}(a_t\mid s_t)}{\pi_{\phi_{old}}(a_t\mid s_t)},
\label{eq:ppo_ratio}
\end{equation}
则策略优化目标为
\begin{equation}
L_{\mathrm{CLIP}}(\phi)=
\mathbb{E}_t\left[
\min\left(
w_t^{\mathrm{IS}}(\phi)\hat{A}_t,
\mathrm{clip}(w_t^{\mathrm{IS}}(\phi),1-\epsilon,1+\epsilon)\hat{A}_t
\right)
\right],
\label{eq:ppo_clip}
\end{equation}
其中 $\hat{A}_t$ 为优势函数估计，$\epsilon$ 为 clip 参数。优势函数由广义优势估计（Generalized Advantage Estimation, GAE）计算：
\begin{equation}
\hat{A}_t=\sum_{n=0}^{\infty}
(\gamma_{\mathrm{RL}}\lambda_{\mathrm{GAE}})^n
e_{t+n}^{\mathrm{TD}},
\quad
e_t^{\mathrm{TD}}=r_t
+\gamma_{\mathrm{RL}}V_{\theta}(s_{t+1})-V_{\theta}(s_t),
\label{eq:gae}
\end{equation}
其中 $\gamma_{\mathrm{RL}}$ 为折扣因子，
$\lambda_{\mathrm{GAE}}$ 为 GAE 参数。外层 RL 决策保持固定步长 $\sigma$，与区间内部的连续物理事件并存。

训练完成后，actor 在每个滚动时刻输出首期每个槽位的请求功率，并通过式\eqref{eq:rl_power_rollout}生成供 MPC 使用的预测域请求功率序列；实际仅执行首期请求功率。critic 的分段线性参数在本轮求解中固定，MPC 根据自身的路径与服务决策构造 $s_{\ell+H}^{\mathrm{MPC}}$，通过式\eqref{eq:pwl_critic_embedding}联合确定终端价值。由此，RL 层同时提供请求功率轨迹和后续运营价值模型，MPC 则优化剩余路径与连续事件安排，形成闭环协同机制；参考推演仅用于生成功率等求解输入，不预先固定本轮终端价值。


\section{Numerical experiments}

% NUMERICAL_EXPERIMENTS_PLACEHOLDER_BEGIN
% 待作者提供实验设置、数据与结果
% NUMERICAL_EXPERIMENTS_PLACEHOLDER_END


\section{Conclusion}

本文针对高速公路换电站网络中预约与随机两类异质用户的运营优化问题，构建了 MPC--RL 分层协同框架。在 MPC 层，本文以滚动时域组织运营：未到预约用户从入口出发、相对日前初始发布路径评价调整，在途预约用户从实时虚拟起点出发、相对最近发布的未执行站序评价调整，站内仍在等待的请求带入下一轮且保留原到站时刻与截止时刻；候选网络按 SOC 分档离线生成、在线裁剪，路径调整成本按实际发布结算一次。站内服务以连续时间刻画：请求在到达后最长等待 $\Delta t$，预约优先、类内按到站顺序服务，槽位电池严格在达到 SOC $1$ 后交付，充电在满电时停止、换入退回电池后立即恢复。在 RL 层，actor 输出每个槽位的请求充电功率作为 MPC 的已知参数，critic 通过分段线性状态价值函数，统一衡量期末库存与未完成请求对后续运营净回报的影响；计划换电路径先于 actor 滚动推演确定，随后在固定功率序列和网络参数下由 MPC 联合优化路径、服务与终端价值，避免请求序列与本轮路径决策之间的循环依赖。

本文的模型边界包括：车辆到达下游节点的预计到达时间为外生给定，未建立随机交通模型；换电视为瞬时事件，未区分服务开始与作业结束；终端价值采用有限维状态编码与分段线性 critic 近似，不能保证完整保留库存分布和多请求依赖，其误差及新增整数变量对求解效率的影响需通过实验评估。后续工作可围绕更精细的交通流建模、换电作业时长以及终端价值的精确刻画展开，并通过系统的数值实验评估框架在不同需求强度、电价模式和库存配置下的表现。


%%----------- 参考文献 -------------------%%
%在reference.bib文件中填写参考文献，此处自动生成

% 当前正文尚无 \cite 命令，添加正式引用后再启用下行参考文献列表。
% \reference

\end{document}
